\section{A limitation of contracting two diameter tensors at the centre}
\label{sec:mixed-lines}

The exact edge functionals on the two diameters give more information
than the wire exponents alone. They also permit a direct test of one
class of mixed identities. The result below excludes every sum formed
by contracting one tensor on each diameter at their shared centre,
when the desired output is the two adjacent wires. It does not test
the opposite-wire crossing used in the main question.

\begin{lemma}[The zero-weight flattening is controlled by edge support]
\label{lem:edge-centre-rank}
Let $T$ be a three-party tensor over any field, and let $r_3(T)$
be its matrix rank across the third party versus the first two.
For the upper support functional with weights $(1/2,1/2,0)$,
\[
 F(T):=\zeta^{(1/2,1/2,0)}(T)\geq\sqrt{r_3(T)}.
\]
\end{lemma}
\begin{proof}
The zero tensor is immediate. In any choice of bases, let $R$ be the
number of distinct pairs of first and second coordinates that occur
in the support. These pairs are precisely the nonzero rows of the
third-party flattening, so $R\geq r_3(T)$. For each pair choose one
supported third coordinate, and give the resulting $R$ support
triples equal probability. The joint entropy of the first two
coordinates is $\log_2 R$. Entropy subadditivity gives
\[
 \tfrac12H(P_1)+\tfrac12H(P_2)
 \geq\tfrac12H(P_{12})=\tfrac12\log_2 R.
\]
For completeness, subadditivity is the nonnegativity of
$\sum_{i,j}p_{ij}\log_2(p_{ij}/(p_i p_j))$.
The maximum over all support distributions is therefore at least
$\tfrac12\log_2 r_3(T)$ in every basis. Taking the minimum over bases
in the support definition, recalled in \cite[equation (3), p.~3]{alp-edge},
and then exponentiating proves the claim.
\end{proof}

\begin{proposition}[Every sum of two-diameter contractions obeys the test]
\label{prop:mixed-diameter-sums}
On the five parties $1,2,3,4,0$, put
$S_q=\langle q\rangle_{12}\langle q\rangle_{34}$, with a
one-dimensional factor at the centre $0$. Suppose $S_q$ is in the
geometric Zariski closure of sums of restrictions of the fixed resources
\[
 T_i(1,3,0)\,U_i(2,4,0),\qquad 1\leq i\leq A.
\]
Here the centre is the zero-weight slot when $F$ is evaluated.
Then
\begin{equation}\label{eq:mixed-diameter-cost}
 \sum_{i=1}^A F(T_i)F(U_i)\geq q^2.
\end{equation}
In particular, let a hypothetical normalized spectral point on the
centred square have diagonal exponent $g>0$, adjacent exponent
$g/\sqrt2$, and restriction $F^g$ on each diameter triple, as in
Corollary~\ref{cor:square-edge-restrictions}. For every such candidate
cover its prescribed values satisfy
\[
 \sum_{i=1}^A\bigl(F(T_i)F(U_i)\bigr)^g
 \geq q^{2g}\geq q^{\sqrt2g}=\psi(S_q).
\]
Thus this class of identities cannot contradict the prescribed values,
for any number of summands or at any permitted scale.
\end{proposition}
\begin{proof}
Let $r_i$ and $s_i$ be the centre flattening ranks of $T_i$ and
$U_i$. After the local maps at the four corners, regard each tensor
as a linear map from its centre factor to the tensor product of its
two endpoint spaces. A centre contraction pairs those two factors.
Consequently the output matrix across $13\mid24$ has rank at most
$\min(r_i,s_i)$. This conclusion also follows by writing the output
as a product of the two centre-flattening matrices and the matrix
of the contracting functional.

Rank subadditivity bounds every sum by
$\sum_i\min(r_i,s_i)$. The locus of matrices of rank at most this
integer is closed, since it is defined by minors, so the same bound
holds in the stated closure. The flattening of $S_q$ across
$13\mid24$ has rank $q^2$. Lemma~\ref{lem:edge-centre-rank} now gives
\[
 q^2\leq\sum_i\min(r_i,s_i)
 \leq\sum_i\sqrt{r_i s_i}
 \leq\sum_iF(T_i)F(U_i),
\]
proving \eqref{eq:mixed-diameter-cost} over the original field and
under geometric base extension alike.

A normalized wire exponent satisfies $0<g\leq1$. For nonnegative
numbers $c_i$ and this range of $g$,
$\sum_i c_i^g\geq(\sum_i c_i)^g$; this follows by applying
$t^g\geq t$ to $t=c_i/\sum_jc_j$, with the zero-sum case immediate.
Use $c_i=F(T_i)F(U_i)$ and \eqref{eq:mixed-diameter-cost}. Finally,
the two adjacent wires each have exponent $g/\sqrt2$, so their
product has value $q^{\sqrt2g}$. This proves the last assertion
directly for the exact sum, without replacing it by the number of
terms times the largest term.
\end{proof}

The opposite-wire crossing
$\langle q\rangle_{13}\langle q\rangle_{24}$ has rank one across
$13\mid24$. The proof therefore gives no useful obstruction for
that output. Identities with other intermediate tensors, additional
line directions, or a different pattern of local operations also
remain outside the proposition.
